\setcounter{subsection}{0}
\section{RESULTADOS}

Este capítulo presenta los resultados de las Fases~3, 4 y~5 del proyecto, ligadas a los objetivos
tercero, cuarto y quinto. Los de los componentes de datos y de recuperación siguen la fase de
evaluación de CRISP-DM, y los del código, el plan del informe de pruebas
(Sección~\ref{sec:calidad}). Cada cifra va con su muestra, su condición y su grado de respaldo: recalculada desde una
salida conservada, respaldada por una salida existente o solo reportada. Cuando dos fuentes
discrepan o falta el programa que produjo una cifra, el texto lo dice en lugar de dar un valor.

% =====================================================================
\subsection{Percepción documental}
\label{sec:percepcion}

El primer criterio del árbol de utilidad exige extraer de una lista escolar, entregada como imagen
o PDF sin estructura, los renglones con su cantidad y su descripción. El 1 de septiembre de 2026,
los prototipos de tesis compararon cinco motores sobre CGS-10, un conjunto de oro de diez documentos
y 1.293 ítems. La métrica es F1 en sus variantes micro, agregada sobre todos los renglones, y macro,
promediada por documento (Tabla~\ref{tab:percepcion}).

\begin{table}[htbp]
\centering
\caption{Percepción documental en los prototipos de tesis (CGS-10, diez documentos)}
\label{tab:percepcion}
\small
\begin{tabular}{@{}lrr@{}}
\toprule
\textbf{Enfoque} & \textbf{F1 micro} & \textbf{F1 macro} \\
\midrule
Gemini (multimodal) & 0,948 & 0,944 \\
Google Vision & 0,673 & 0,672 \\
GOT-OCR2.0 & 0,597 & 0,569 \\
Docling & 0,364 & 0,424 \\
OCR con reglas heurísticas & 0,298 & 0,419 \\
\bottomrule
\end{tabular}
\end{table}

Los valores se recalcularon desde la salida conservada de la comparación, con su regla de
emparejamiento (similitud mínima de 0,75), sin volver a ejecutar los motores. El acta del OE3, que
está pendiente de firma y no forma parte de la entrega, consigna 0,925, un valor que no proviene de
esa salida. La Figura~\ref{fig:bench-F3-2} añade la latencia por documento que registra la misma
salida; puede incluir esperas por el límite de tasa, de modo que no sirve para verificar DP-04 del
SRS. La latencia de la percepción en el sistema construido sigue pendiente: la verificación usó una
percepción simulada, y la medición con el modelo real corresponde al segundo nivel de las pruebas
(ESC-D1 del ATAM).

La comparación tiene tres amenazas a su validez. La principal es el oro: CGS-10 se armó corrigiendo,
en cada documento, la extracción de uno de los motores evaluados, sin un segundo anotador. Coincide
con la salida de Gemini en el 86,7\,\% de los ítems literales y en el 93,1\,\% sin tildes ni
puntuación, así que el 0,948 no es independiente de esa salida. La segunda es la muestra: diez
documentos y una sola corrida por motor. La tercera es que el sistema construido usa otro modelo de
Gemini, con temperatura 0 y sin semilla. La ventaja del modelo multimodal frente a las reglas es
amplia, pero su magnitud está sesgada a favor del motor elegido.

\input{\graficasdir/F3-2.tex}

% =====================================================================
\subsection{Recuperación y normalización}
\label{sec:recuperacion}

El 2 de septiembre, los prototipos de tesis compararon cuatro estrategias sobre PostgreSQL con
pgvector: léxica pura, léxica expandida, vectorial densa e híbrida con fusión recíproca de rangos
(RRF, $k=60$). El sistema construido adoptó en su lugar la fusión aditiva de una sola fase
descrita en la Sección~\ref{sec:normalizacion}.

\subsubsection{Latencia}

La léxica pura tuvo un p50 de 6,66\,ms y una media de 40,60\,ms; la léxica expandida, de 6,59 y
52,76\,ms; la vectorial densa, de 46,51 y 106,61\,ms, y la híbrida con RRF, de 53,17 y 155,51\,ms.
La medición registra el tiempo y el número de resultados, no la relevancia, y no documenta un
calentamiento; la media está entre dos y ocho veces el p50 por las colas largas. En la misma fecha, con 50 consultas
en CPU, pgvector con índice HNSW tuvo una media de 5,06\,ms (p95 de 6,22\,ms); la búsqueda exacta en
pgvector, 363,68\,ms, y el índice HNSW de FAISS, 0,098\,ms. El QA acepta los 5,06\,ms de pgvector
por sus garantías transaccionales. La latencia de la recuperación del sistema construido sigue
pendiente; la prueba de carga de la verificación midió la consulta del catálogo, no la recuperación
de candidatos (Sección~\ref{sec:calidad}).

\subsubsection{Calidad de la recuperación}

Con la biblioteca de MTEB \citep{muennighoff_mteb_2023} se compararon dos modelos de embedding en
la tarea SpanishPassageRetrievalS2P, de pasajes de salud y de dominio distinto al escolar.
multilingual-e5-large obtuvo un MRR@10 de 0,6163, frente a 0,4735 de
paraphrase-multilingual-mpnet-base-v2; las salidas no registran la versión de la biblioteca. Los
550,14\,ms y 69,88\,ms que algunos documentos técnicos presentan como latencia por consulta resultan
de dividir el tiempo total de la tarea entre un número aproximado de consultas (1.088,20\,s entre
1.978,03 en el modelo grande), así que no lo son. El resultado tampoco describe el sistema
construido, que calcula los embeddings con gemini-embedding-001 desde el 23 de septiembre de 2026.

Sobre Recall@10 hay una discrepancia sin resolver. Un informe de los prototipos reporta 1,0000. La
única salida conservada, la de las 50 consultas, da 0,10 tanto para FAISS como para pgvector con
HNSW frente a la búsqueda exacta, y el programa que la produjo no se conserva. En el banco de
alternativas, la búsqueda exacta sobre 10.000 vectores sintéticos mantuvo un Recall@10 de 1,0. Este
capítulo no afirma un valor de Recall@10 para los índices aproximados.

\input{\graficasdir/F5-2.tex}

\subsubsection{Acuerdo en la clasificación del corpus}

El 7 de septiembre, los prototipos de tesis midieron un kappa de Cohen de 0,7331 sobre 200 pares de
12 documentos, que se reprodujo desde las entradas conservadas; su intervalo de confianza del
95\,\% por remuestreo es [0,670; 0,794]. Un segundo valor, 0,7049 con 25 documentos, no se puede
recalcular. El kappa no mide el acuerdo entre dos personas: compara la categoría que asigna una
heurística con la que asigna un modelo de lenguaje, en 45 de las 48 categorías del corpus. Por eso
no cumple la meta del SPMP, un kappa de 0,7 o más con un segundo anotador humano sobre al menos 30
documentos, y el límite inferior del intervalo queda además por debajo de 0,7.

\subsubsection{Umbral de resolución automática}

Los tres valores del umbral (Sección~\ref{sec:normalizacion}) operan sobre escalas distintas y no
son intercambiables. El de 0,630 se calibró con 542 ítems y un AUC de 0,880, pero 186 de sus
etiquetas (34,3\,\%) se propagaron a descripciones idénticas. La muestra efectiva es de 356
descripciones, así que los intervalos que tratan los 542 ítems como independientes son más estrechos
de lo justificado, y el mismo conjunto sirvió para elegir el umbral y para medirlo. La
Figura~\ref{fig:bench-F5-2} muestra cómo cambian la cobertura y la precisión del primer candidato
al desplazarlo; en 0,630 se resuelven 354 de los 542 ítems (65,3\,\%), y 307 de ellos (86,7\,\%)
correctamente.

La única medición del umbral de 0,75 pertenece a una línea de desarrollo del 4 de octubre, anterior
a la versión evaluada y no integrada en el sistema construido. Un conjunto de oro de 2.588
renglones, tomados de 4.719 ítems y 98 conceptos, midió la regla de una versión anterior. Las
versiones evaluada y construida conservan el umbral, pero cambiaron el análisis del renglón y el
catálogo, de modo que la medición no las describe. Aquella regla resolvió de forma directa 2.252 renglones, 2.234 de ellos
correctos (precisión directa de 0,992), dejó 298 en la zona gris (11,51\,\%) y 38 fuera de la
resolución. Sobre todas las líneas, el acierto fue de 0,9784; si en la zona gris decide el
acudiente, de 0,9737, y si decide el sistema, de 0,8632. La medición solo ejercita el canal léxico,
porque los 609 productos tenían el embedding vacío, y 1.138 de sus renglones son coincidencias
exactas triviales. El umbral de 0,75 sigue sin calibrar.

\subsubsection{Extremo a extremo}

La cifra de 83\,\% sobre 60 sublistas y 2.176 ítems, del servicio de búsqueda anterior, mide la
fracción de ítems con al menos un candidato recuperado (83,04\,\%), no la resolución. En esa
corrida, 769 ítems (35,3\,\%) se resolvieron de forma automática, 795 (36,5\,\%) pasaron a revisión
y 612 (28,1\,\%) quedaron sin correspondencia. Se ejecutó sin conexión, en una base de prueba, sobre
un espejo de 838 productos con los 48 nombres más cortos de cada categoría.

De las primeras versiones del sistema construido hay una medición de distribución, no de precisión. Entre el 21 y el 23 de
septiembre, un ensayo de huecos aplicó el pipeline a 14 PDF de listas, con 199 páginas y 4.414
renglones. Se reconocieron 2.202; 926 terminaron en abstención (20,98\,\%), 626 sin coincidencia,
66 como hueco del catálogo y 594 en una clase «fuera» que los documentos no definen. El ensayo dejó
1.397 faltantes únicos para la curaduría del catálogo.

\subsubsection{Auditoría de los kits}

Las mediciones anteriores evalúan la resolución de cada renglón. Una línea de afinamiento de los
kits evaluó además el resultado que recibe la familia. Una auditoría previa, del 7 de octubre,
revisó los kits de un grado, Séptimo: al incorporar los atributos del catálogo, las líneas del
Ahorrador que cumplen pasaron de 31 a 49 y las que no cumplen, de 30 a 3, y de ahí salieron los
atributos y el filtro por renglón. El 11 y el 12 de octubre de 2026 se auditaron cuatro veces, renglón por renglón, los kits
Ahorrador y Premium de las trece listas de 2026 del Colegio Agustiniano Ciudad Salitre. Un renglón
cumple si los dos perfiles cobran lo que pide la lista, y son comprables todos los renglones salvo
los que Tukit no vende, como los uniformes o los textos con título. Entre una auditoría y la
siguiente se ajustaron la comparación de atributos y la curaduría de las listas y del catálogo, y
antes de la cuarta el catálogo pasó de 679 a 1.030 productos con precio (Tabla~\ref{tab:kits}).

\begin{table}[htbp]
\centering
\caption[Auditorías de los kits del Colegio Agustiniano Ciudad Salitre]{Auditorías de los kits del Colegio Agustiniano Ciudad Salitre, 11 y 12 de octubre de
2026; las dos últimas, posteriores a la versión evaluada (cifras reportadas; costo de los trece kits de cada perfil)}
\label{tab:kits}
\small
\begin{tabular}{@{}lrrrrr@{}}
\toprule
\textbf{Auditoría} & \textbf{Comprables} & \textbf{Huecos} & \textbf{Cobros} &
\multicolumn{2}{c}{\textbf{Costo (millones de pesos)}} \\
\cmidrule(l){5-6}
 & \textbf{que cumplen} & \textbf{declarados} & \textbf{sin cumplir} & \textbf{Ahorrador} &
\textbf{Premium} \\
\midrule
1 (11 oct.) & 263 de 888 (29,6\,\%) & 240 & 589 & 13,9 & 32,0 \\
2 (11 oct.) & 465 de 888 (52,4\,\%) & 385 & 242 & 9,9 & 17,6 \\
3 (11 oct.) & 617 de 886 (69,6\,\%) & 220 & 49 & 9,0 & 14,2 \\
4 (12 oct.) & 758 de 885 (85,6\,\%) & 69 & 58 & 8,1 & 17,3 \\
\bottomrule
\end{tabular}
\end{table}

El cumplimiento subió en cada auditoría. Los cobros de un producto que no corresponde a lo pedido
bajaron de 589 a 49 en las tres primeras y subieron a 58 en la cuarta, sobre todo por 23 renglones
en los que el Ahorrador cobra un cuaderno que no dice ser grande. Los huecos no bajaron de forma
continua: los 240 iniciales incluían 180 renglones que Tukit no vende, y de los 69 finales, 45 son
de electrónica, que las tiendas registradas no ofrecen. El Premium se encareció en la cuarta
auditoría porque cobra 151 renglones que antes eran huecos. Las cifras provienen de los informes de
auditoría y cubren un solo colegio, cuyas listas se cargaron desde datos curados. Dependen de una
clasificación humana de cada renglón. Las auditorías 1 y 2 midieron versiones del 5 y el 8 de
octubre, que forman parte de la historia de la versión evaluada, y la segunda tiene su mismo código
de producto. Las auditorías 3 y 4 midieron versiones intermedias del afinamiento, que el sistema
construido integra con cambios posteriores en la tolerancia de las medidas, el tamaño del cuaderno
y el catálogo, de modo que ninguna auditoría mide el sistema construido. El anexo
CRISP-DM detalla cada iteración y sus límites.

% =====================================================================
\subsection{Optimización}
\label{sec:optimizacion}

Sobre la calidad del optimizador (Sección~\ref{sec:optimizacion-dev}) hay tres evidencias de
alcance distinto. El 31 de julio, el banco de alternativas comparó NSGA-II con un algoritmo voraz y
uno aleatorio en nueve instancias sintéticas: halló el óptimo en ocho y quedó a un 0,07\,\% de él,
frente al 11,92\,\% del voraz y el 159,34\,\% del aleatorio. Solo se conserva el informe, y el
problema tenía tres objetivos, de modo que no valida el optimizador del sistema construido. En este,
las pruebas calculan un hipervolumen \citep{zitzler_multiobjective_1999} en dos dimensiones sobre un
ejemplo de enumeración y comprueban, con 65.536 combinaciones, que el frente de NSGA-II no tenga
soluciones dominadas y que la semilla lo haga reproducible; no hay registro de su ejecución. Un
banco del 4 de octubre, que forma parte de la versión evaluada, compara NSGA-II con un óptimo exacto
por programación dinámica sobre 174 listas del corpus. En el Ahorrador lo iguala en el 100\,\% de
las listas. Con presupuesto, el kit recomendado daba el 40\,\% de la calidad posible; después del
ADR 0020 cabe siempre en el tope e iguala el óptimo. Son cifras reportadas con una sola semilla y no
reproducidas, y la resolución usó solo el canal léxico y una foto fija del catálogo. El hipervolumen
sigue sin calcularse fuera de la prueba automatizada.

Sobre 29 listas reales, otra medición reportada dio 5,3~s de media, 6,3~s de mediana y 9,0~s de
máximo para producir todos los kits de una lista, en una CPU de 16 núcleos sin otras cargas. El
máximo supera los 8~s de DP-03. No es la prueba del escenario ESC-D4 del ATAM, que pide al menos 30
ejecuciones y el percentil 95, y midió el optimizador anterior a las correcciones del 4 y el 7 de
octubre. El tiempo en el peor caso de la versión evaluada sigue pendiente.

% =====================================================================
\subsection{Verificación de la versión evaluada}
\label{sec:calidad}

Los porcentajes de cobertura que citan algunos documentos sobre los microservicios anteriores no
son mediciones. Son los umbrales mínimos que su integración continua exigía para aprobar: 90\,\% en
cuentas y en comercio, 84\,\% en instituciones y 80\,\% en búsqueda. Esas ejecuciones terminaron con
éxito el 18 de septiembre de 2026, de modo que la cobertura alcanzó al menos el umbral, pero no se
conserva el valor medido.

La versión evaluada se verificó con un plan de diez fases, distintas de las cinco del proyecto:
la línea base (0), la matriz de requisitos no funcionales (1), las pruebas unitarias, de
integración, funcionales con conjuntos de oro, adversariales de la IA y de los patrones de
resiliencia (2 a 6), la degradación de componentes (7), la carga (8), la telemetría (9) y el
informe (10). Es un resultado de la Fase~4 de la VFP, que espera pruebas automatizadas del motor de
optimización y de las interfaces, y cubre en parte las pruebas funcionales, no funcionales y de
carga de la Fase~5. El informe de pruebas de la entrega reúne el método, las evidencias y los
resultados. Las fases 0 a 6 corrieron sobre la versión de integración del 8 de octubre de 2026,
cuyo código de producto es el mismo de la versión evaluada salvo la integración continua y dos
valores por omisión de la base de datos. El primer nivel de las fases 7 a 9 corrió sobre la versión
evaluada, en una sola corrida, en una máquina virtual de 4 procesadores y 8~GiB y con la percepción
simulada. La verificación no hizo ninguna llamada real al modelo: el simulador atendió 2.342
percepciones.

Con las pruebas propias del sistema pasan 1.647 pruebas, y la cobertura es de 72,43\,\% de 12.039
sentencias. Con las que añadió la verificación, pendientes de integrar, pasan 3.911, otras 129 son
fallas esperadas que documentan defectos y la cobertura sube a 94,32\,\%. La integración continua
ejecuta las pruebas, pero no mide la cobertura, y su estado en la versión evaluada no se comprobó.
El modelo del catálogo difiere de sus migraciones en dos valores por omisión. De las 60 filas de la
matriz que liga cada meta con su ensayo, se ejecutaron 36: 7 cumplen, 7 cumplen en lo medido a la
espera de los niveles 2 y 3, y 22 no cumplen. Otras 13 están pendientes, 2 esperan un hallazgo ya
anticipado, 8 no se pueden medir con los componentes actuales y 1 es el protocolo con usuarios.
Quedan 78 hallazgos abiertos: 14 de severidad alta, 56 de severidad media y 8 de severidad baja.

El primer nivel de las fases 7 a 9 dejó 18 ensayos, de los que 14 cumplen y 4 no. En la carga, una
rampa hasta 100 usuarios generó 27.894 peticiones sin errores, y el percentil 95 de la consulta del
catálogo fue de 22~ms con un usuario y de 160~ms con cien, por debajo de los 300~ms de la meta. No
hubo ningún rechazo por límite de tasa porque la versión evaluada no tenía limitador, de modo que en
ella esa meta no se cumple; el sistema construido añade uno con Redis, pendiente de medición. En la degradación, siete de diez caídas de un componente cumplieron: los demás servicios
siguieron respondiendo sin errores y se recuperaron en 2,5~s o menos. Con el catálogo colgado, el
disyuntor abrió a los 26,3~s, dentro de los 30~s del escenario ESC-R2, porque cada intento falló
en unos 5~s. Si cada intento agotara su espera, el modelo da 300,8~s desde inteligencia y 150,8~s
desde comercio. Ese peor caso no se midió, y el hallazgo queda acotado, no corregido. Pasados los
30~s, con diez llamadas simultáneas, el disyuntor deja pasar las diez en lugar de una sola de
prueba.

Las tres caídas que no cumplieron dejaron ver tres limitaciones de la versión evaluada. Con el
catálogo detenido, subir una lista respondía con un error 500 y la lista no se guardaba. Después de
caer la base de datos o el intermediario de mensajes, los procesos de fondo morían y no se
reiniciaban. Y un mismo pago producía dos avisos al acudiente, sin cobro doble. El sistema
construido corrige las tres. Con el catálogo caído, la subida responde 202, guarda lo percibido y
la lista se importa cuando el catálogo vuelve (ADR 0031). Los procesos de fondo sobreviven a las
caídas pasajeras de la base o del intermediario, aunque un error de programación todavía los
detiene (ADR 0032), y un evento produce un solo aviso (ADR 0033). La corrección se verificó el 20
de octubre de 2026 en una versión de trabajo, antes de integrarla: la subida responde 202 en
33~ms y se recupera en 2,7~s, los procesos de fondo se recuperan en 7,1~s y 6,9~s, y el pago
produce un solo aviso. Esa verificación no volvió a correr los demás ensayos, y la corrida
completa sobre el sistema construido está pendiente.

Quedan pendientes la corrida completa sobre el sistema construido y la medición de lo que añadió
después de la versión evaluada: las colas por RabbitMQ, la caché y el límite de tasa con Redis, las sugerencias de dirección, el alta de
colegios no aliados y la comparación por atributo. También lo están los niveles 2 y 3 de la
degradación, la carga y la telemetría, con la percepción real, el pico estacional y las caídas
combinadas, y el estudio con usuarios de la Fase~5
(Sección~\ref{sec:oe5}). En la telemetría, las trazas no llevan el identificador de correlación, de
modo que una carga no se puede seguir de extremo a extremo.

% =====================================================================
\subsection{Auditoría de seguridad y privacidad}
\label{sec:defectos}

Dentro del OE2 se hicieron dos auditorías, listadas en la Tabla~\ref{tab:anexos}: la de
autorización y la de tratamiento de datos personales (Anexo~M), esta conforme a la Ley 1581 de 2012
\citep{congreso_ley1581_2012} y al Decreto 1377 de 2013 \citep{presidencia_decreto1377_2013}. Las
dos revisaron el código de los microservicios anteriores. La de autorización revisó 152 puntos de
entrada y encontró 22 sin una comparación completa de ámbito: 8 sin ninguna y 14 con una parcial.
El Anexo~M registra catorce defectos, dos críticos y cinco de severidad alta.

\subsubsection{Defectos de los microservicios anteriores}

El de mayor alcance estaba en las políticas de seguridad a nivel de fila de las tablas con datos de
estudiantes. La condición base concedía visibilidad de la tabla completa cuando el contexto de
usuario no estaba establecido, de modo que el error más probable en un sistema distribuido, que un
servicio olvide fijar ese contexto, otorgaba acceso irrestricto. Una segunda cláusula daba acceso
completo al rol de administrador comparando solo el identificador del usuario, sin ancla de ámbito,
y ese defecto motivó el modelo de autorización con ámbito del capítulo de diseño. El ADR 0015 lo
registra en 16 tablas de cinco servicios y advierte que corregirlo va antes que su propia decisión;
las migraciones posteriores conservaban la apertura, sin pruebas de regresión. Entre los demás, el
manejador genérico de errores de validación devolvía al cliente el valor rechazado, contraseña
incluida, y la supresión de datos se marcaba como completada sin confirmación de los servicios.

\subsubsection{Estado en el sistema construido}

El sistema construido no tiene políticas de seguridad a nivel de fila. El ámbito se controla en la
aplicación, con la cadena de contención de identidad, que falla cerrada
(Sección~\ref{sec:autorizacion}); las cifras de 152 y 22 son las de la auditoría, no una medición
propia. De los catorce defectos del Anexo~M, dos están corregidos: el valor rechazado ya no vuelve
al cliente, y el correo del acudiente ya no queda en un registro que la supresión no alcanza. El
texto de la lista sin depurar se corrigió para la lista oficial, pero sigue en el documento que
carga el acudiente; la depuración del nombre se sustituyó por la minimización
(Sección~\ref{sec:percepcion-dev}), salvo en ese mismo documento. Los tres defectos de la supresión
cambiaron de forma, pero la supresión sigue sin mecanismo: el estado de propagación existe en el
modelo, sin publicador ni consumidor del evento.

Cuatro siguen vigentes: la gestión de cuentas no deja registro de auditoría, el consentimiento por
finalidad no se escribe ni se consulta, la respuesta de la importación de listas guarda el extracto
completo del documento y ninguna tabla de auditoría tiene una política de purga. Otro no tiene
efecto mientras no se escriba el registro de auditoría. Uno es parcial: las preferencias y los pesos
del acudiente están cifrados, pero la razón de la recomendación queda en claro y no se borra. El de
la anonimización no es evaluable, porque no hay ese proceso.

Un defecto de comercio quedó corregido en la versión evaluada. La consulta de los pedidos del
acudiente toma el usuario de la sesión, validada con identidad, e ignora el que llegue en la
consulta; el detalle, el pago y la cancelación de un pedido propio también exigen sesión, y un
pedido ajeno responde 404. La puerta de enlace sigue sin exigir sesión.

En el sistema construido siguen abiertos otros defectos. La consulta de contactos del servicio de identidad devuelve el correo de
cualquier usuario activo sin autenticación, y las rutas de los paneles de catálogo, inteligencia y
operación no exigen sesión ni rol. Cuatro caminos crean matrículas, y solo uno, el registro de un
hijo desde la cuenta del acudiente, escribe el consentimiento del menor. El celular del acudiente se
guarda sin cifrar en la base de identidad.

Las pruebas propias del sistema no detectaron ninguno de estos defectos. Medir cuánto código
ejecutan no dice si buscan las condiciones adversas correctas, y varias comprueban garantías leyendo
el código fuente en lugar de ejecutar el comportamiento. La verificación de la versión evaluada
(Sección~\ref{sec:calidad}) los confirmó y encontró otros. Un kit se puede comprar sin sesión, aun
mientras espera revisión, y varias rutas internas aceptan llamadas sin verificar quién llama. Las
cinco bases de datos comparten un mismo rol de acceso, de modo que las credenciales de cualquier
servicio leen los usuarios de identidad. La página del documento cargado sale sin depurar hacia el
proveedor del modelo (\mbox{PE-06}), y el nombre ficticio de las pruebas quedó guardado. El retracto existe
en comercio, pero ninguna aplicación lo ofrece y la puerta de enlace no lo expone; además, se acepta
pasado el plazo legal y no libera la reserva. Un pedido de 100.000 unidades fue aceptado para pago.
La respuesta de la puerta de enlace expone direcciones internas, el mensaje de un trabajo fallido
contiene la consulta SQL con sus parámetros y, sin la llave de cifrado, el cifrador guarda en claro
sin avisar.

% =====================================================================
\subsection{Validación con usuarios (OE5): planificada y pendiente}
\label{sec:oe5}

El quinto objetivo específico de la VFP pide evaluar el sistema con pruebas funcionales, no
funcionales, de carga y de usabilidad con usuarios reales, y medir la satisfacción y la reducción de
la carga cognitiva frente al proceso tradicional de compra. La verificación de la
Sección~\ref{sec:calidad} cubre en parte las pruebas funcionales, no funcionales y de carga. Las
sesiones con usuarios, con SUS y NPS, son una actividad planificada de la Fase~5 que todavía no se
ha hecho, y sus resultados no se estiman aquí. Los
escenarios de prueba del ATAM fijan los criterios de aceptación, entre ellos un SUS de 68 o más y 50
usuarios concurrentes con una tasa de error inferior al 1\,\%.

El método es el de la VFP. Antes de involucrar participantes exige un protocolo ético aprobado
\citep{lazar_research_2017}, conforme a la Ley 1581 de 2012 y al Decreto 1377 de 2013, porque la
interacción toca datos de menores. Prevé sesiones observadas con tareas representativas, el
cuestionario SUS \citep{brooke_sus_1996,bangor_empirical_2008,lewis_factor_2009} con al menos cinco
participantes e idealmente veinte o más, verbalización concurrente, el índice NPS
\citep{reichheld_one_2003} y, para la carga cognitiva, el conteo de interacciones discretas (clics,
búsquedas y validaciones manuales) frente al flujo tradicional. UI-07 del SRS v6.1, el QA, la WBS y
el ATAM usan NASA-TLX; como no hay una decisión registrada de ese cambio, el OE5 se juzga con el
método de la VFP.

El informe de pruebas propone un protocolo para esas sesiones. Recluta 30 acudientes que hayan
comprado útiles escolares en los dos últimos años y no hayan usado Tukit, para llegar a 26 sesiones
válidas; si no se alcanzan, el informe dará la muestra lograda y la potencia real, sin declarar
cumplimiento. Cada sesión, de unos 45 minutos, compara la compra en dos o tres tiendas en línea con
la obtención del kit en Tukit hasta un pago simulado, en orden contrabalanceado, con un estudiante
ficticio y sin datos de menores. Quedan cinco decisiones abiertas, entre ellas quién facilita y
observa las sesiones y quién tramita el aval de ética, si la universidad lo exige. El protocolo
ético aprobado, los consentimientos y el análisis no forman parte de esta versión.
